\section{Introduction}
\label{sec:intro}

Deep research systems automate the process of researching a question. A language model issues web searches over many rounds, reads the retrieved pages, decides after each round whether more searching is needed, and finally writes a report in which each statement cites the page it came from. OpenAI Deep Research \citep{openai2025deepresearch}, released in February 2025, and STORM \citep{shao2024storm}, introduced by the Stanford OVAL lab in 2024, are two of the most widely used systems of this kind.

Both systems are designed to produce surveys. Asked whether the Federal Reserve should keep interest rates high, they present the evidence on each side and leave the judgment to the reader. A survey is useful for preliminary research, but many professional tasks end in a decision. A lawyer drafting a brief argues for one side of a case, and an analyst advising on an investment recommends whether to buy or sell. In both settings the reader expects a clear position, a small number of claims that support it, and evidence behind each claim, which a balanced survey does not provide.

The difference between a survey and an argument also suggests a different way to organize research. A survey can grow from the bottom up, because each search can suggest the next and the report describes whatever the searches find. An argument, by contrast, has a structure that is known before any search begins. The main claim depends on a few supporting claims, and each supporting claim depends on evidence that can be retrieved. Research for an argument can therefore be planned in advance, by decomposing the question into sub-questions that are each narrow enough to answer by search and then combining the answers in the order the argument requires.

\sys{} implements this plan. Its name joins the two parts of the method, deep research and decomposition. Given a question, a language model first generates a directed acyclic graph (DAG) of sub-questions, in which every internal node carries instructions for combining the answers of its children. The leaves are then answered by retrieval-augmented generation, and each parent is composed from its children in topological order. When a composed answer is incomplete, the pipeline adds a sub-tree of new questions beneath it. Finally, the answered graph is turned into an outline and a report that argues for the answer at the root. \Cref{fig:pipeline} shows the three phases.

The pipeline was compared with OpenAI Deep Research and STORM in two evaluations, six automated debates and ten reports scored blind by one rater. It received the highest scores for argument quality and took a firm stance far more often than either baseline. A further analysis of the retained reports, which does not rely on a rater, supports the stance result but qualifies two others. Its reports cite densely but draw on few sources, and their conclusions are often more qualified than their titles. An automatic test of citation attribution then shows that the citations in the evaluated reports rarely point to the source a claim came from. The cause lies in how citation numbers were handled as answers were combined, and it applies to any system that composes cited text in stages.

\Cref{sec:problem} defines the task and the requirement that citations survive composition. \Cref{sec:method} describes the pipeline, \Cref{sec:experiments,sec:results} present the original evaluation and its limits, and \Cref{sec:analysis} analyzes the content of the reports. \Cref{sec:citations} presents the attribution test and traces the citation failure to its source, and \Cref{sec:discussion} states what the results do and do not support.
